\documentclass[12pt]{article}
\usepackage[utf8]{inputenc}
\usepackage[T1]{fontenc}
\usepackage{hyperref}
\title{Seberuang Language Resources for K-12 Curriculum Development}
\author{Compiled by Claude Sonnet 5.5 (claude-sonnet-5-5) for Andr\'as Kornai}
\date{}
\begin{document}
\maketitle

\section{Introduction}
Seberuang (ISO 639-3 sbx) is an Austronesian language of West Kalimantan, Indonesia. It is spoken by the Seberuang Dayak, mainly in Sintang Regency (Sepauk, Tempunak and parts of Sintang district) and in Belitang district of Sekadau Regency. Wikipedia cites about 37,000 speakers (2007) and mostly classes it as Ibanic. Some sources call it Malayic. The searches found no school materials, primers or dictionaries. This dossier is therefore short.

\section{Primary Educational Resources}
No primers, textbooks or local-content (muatan lokal) curricula for Seberuang turned up in the searches, including Indonesian-language searches. No evidence was found that it is taught in schools.

\section{Linguistic and Reference Materials}
\begin{itemize}
\item Wikipedia, Seberuang language (overview): \url{https://en.wikipedia.org/wiki/Seberuang_language}
\item ``Dialek sosial bahasa Dayak Seberuang di Desa Sukau Bersatu, Kecamatan Sepauk, Kabupaten Sintang'' (journal article, KANSASI): \url{https://jurnal.stkippersada.ac.id/jurnal/index.php/KAN/article/view/949}
\item Article on reduplication in a Seberuang dialect (Sekubang, Sepauk), journal Sintesis: \url{https://e-journal.usd.ac.id/index.php/sintesis/article/download/922/729}
\item Numeral system data page (Max Planck, Alexander D. Smith): \url{https://lingweb.eva.mpg.de/channumerals/Seberuang.htm}
\item Article on healing mantras of Seberuang practitioners (SEBASA journal): \url{https://e-journal.hamzanwadi.ac.id/index.php/sbs/article/download/33156/8172/206525}
\end{itemize}

\section{Religious and Community Resources}
Joshua Project has a people-group profile for the Seberuang: \url{https://joshuaproject.net/index.php/people_groups/14806}. No Seberuang scripture translation or community organisation was found.

\section{Digital and Audio Resources}
Global Recordings Network lists a page for code sbx: \url{https://glorecnet-website.k.globalrecordings.net/en/language/sbx}. I did not check what recordings it holds.

\section{Curriculum Development Context}
Little published material exists. A curriculum would have to start from new documentation: an orthography, a word list and recorded texts, all done with the community. The dialect, reduplication and mantra studies above are the only linguistic starting points found. The relationship to Malay and Indonesian, and the classification question (Ibanic versus Malayic), should be settled before materials are written.
\end{document}